\section{Recursive Drift and Contextual Entropy}
\label{sec:recursive-drift}

Recursive AI-assisted software systems do not typically fail through one catastrophic decision; they degrade through repeated, locally reasonable steps that gradually separate intent from execution. We refer to this phenomenon as \textit{recursive drift}: the progressive divergence between system intent, generated artifacts, execution context, and human understanding across recursive execution cycles.

Recursive drift is gradual, compounding, and difficult to detect early. Each delegation cycle introduces abstraction, context compression, and interpretation uncertainty. These effects are manageable in isolation, but recursive continuation converts small discrepancies into structural misalignment over time. As in cybernetic control systems, stability depends on timely, high-fidelity feedback loops rather than throughput alone \cite{Wiener1948Cybernetics,Beer1972BrainFirm}.

\subsection{Sources and Categories of Drift}

In recursive software workflows, drift emerges from multiple interacting categories:
\begin{itemize}
  \item \textbf{Contextual divergence:} Context windows, summaries, and prompt rewrites compress prior state, often dropping assumptions that were locally obvious but globally important.
  \item \textbf{Architectural inconsistency:} Delegated agents evolve artifacts under different local heuristics, producing naming, interface, and modularity inconsistencies that violate shared structure \cite{Parnas1972Modularization}.
  \item \textbf{Execution misalignment:} Planned intent and realized implementation diverge when acceptance criteria are interpreted differently across recursive handoffs.
  \item \textbf{Reasoning fragmentation:} Decision rationale becomes distributed across commits, prompts, reviews, and generated outputs, weakening traceability and collective understanding.
  \item \textbf{Temporal desynchronization:} Agents and humans operate over stale assumptions when dependencies or policies change between cycles; asynchronous updates create ordering hazards analogous to distributed systems coordination \cite{Lamport1978TimeClocks,Vogels2009EventuallyConsistent}.
\end{itemize}

\begin{table}[H]
  \centering
  \small
  \caption{Representative recursive drift sources, their effects, and practical mitigation anchors.}
  \label{tab:drift-sources}
  \begin{tabularx}{\linewidth}{l X X}
    \toprule
    \textbf{Source} & \textbf{Primary effect} & \textbf{Mitigation anchor} \\
    \midrule
    Context compression & Omitted assumptions propagate into later plans and reviews. & Synchronization checkpoints that reconstruct scope from canonical artifacts. \\
    Architectural inconsistency & Interfaces and naming drift across delegated implementation units. & Architecture-focused audit passes with explicit compatibility review. \\
    Execution misalignment & Delivered outputs satisfy local interpretations but miss system intent. & Task contracts with concrete acceptance criteria and review gates. \\
    Temporal desynchronization & Work continues on stale constraints after policy or dependency change. & Milestone-level state refresh before further recursive continuation. \\
    \bottomrule
  \end{tabularx}
\end{table}

\subsection{Compounding Recursive Amplification}

Recursive systems amplify small deviations because each cycle conditions the next. For instance, a minor naming mismatch can become interface drift. A missing rationale can become policy non-compliance, and a stale summary can become a faulty execution plan. Once artifacts are reused as context for future generation, local errors are recursively reinterpreted as constraints.

This amplification pressure is intensified by autonomous delegation chains, where AI systems produce tasks, plans, and artifacts for subsequent AI systems. The result is not merely isolated defects but \textit{drift accumulation}: a trajectory in which contextual divergence, artifact inconsistency, and synchronization overhead rise together.

\begin{figure}[H]
  \centering
  \includegraphics[width=\linewidth]{figure2.png}
  \caption{Recursive drift accumulation: locally small divergences between intent, artifacts, and execution context compound into system-level misalignment and rising synchronization pressure across recursive execution cycles.}
  \label{fig:figure2}
\end{figure}

\Cref{fig:figure2} illustrates how locally small divergences compound into system-level misalignment and rising synchronization pressure across recursive execution cycles.

\subsection{Human Cognitive Limits and Synchronization Pressure}

Recursive drift is both a technical and human coordination problem. Human operators work under bounded rationality and finite attention, so they cannot continuously reconstruct full system state from distributed, fast-moving artifacts \cite{Simon1957ModelsMan,Sweller1988CognitiveLoad}. In mixed human-AI workflows, cognition is distributed across tools and participants. As a result, misalignment often emerges at handoff boundaries rather than at any single decision point \cite{Hutchins1995CognitionWild}.

As recursion deepens, the cost of maintaining shared understanding increases faster than local execution speed. This creates \textit{synchronization pressure}: the growing need for explicit checkpoints that reconcile intent, artifacts, and decisions before further delegation.

\subsection{From Drift Detection to Reflective Auditing}

Because recursive drift compounds and remains partially latent until late-stage failures, reactive debugging alone is insufficient.

Recursive AI-assisted systems require reflective synchronization mechanisms that periodically audit assumptions, surface divergence early, and bound continuation when alignment confidence drops. This motivates the next section, which develops reflective auditing as the primary control process for drift detection and alignment recovery (\Cref{sec:reflective-auditing}).
